\section{Hypotheses}
\label{sec:hyp}

The research questions are operationalised as two testable statistical hypotheses.

\textbf{Hypothesis 1 --- GLOF inflation:} GLOF-inclusive annual maxima from the 2011--2018 period
are statistically incompatible with the pre-GLOF reference GEV distribution.
\begin{itemize}
  \item[$H_0^{(1)}$:] The 2011--2018 full annual maxima belong to the same population as
        the pre-GLOF reference period (no distributional shift).
  \item[$H_1^{(1)}$:] The 2011--2018 full annual maxima are drawn from a stochastically
        larger distribution (GLOF-inflated tail).
\end{itemize}
Rejection of $H_0^{(1)}$ at $\alpha = 0.05$ confirms that the GLOF signal inflates the
tail of the annual maximum series.

\textbf{Hypothesis 2 --- Regime stationarity:} Once the direct GLOF contribution is removed by
masking, the 2011--2018 annual maxima are statistically indistinguishable from the
pre-GLOF reference, indicating that the underlying natural-flood regime remained
stationary.
\begin{itemize}
  \item[$H_0^{(2)}$:] The 2011--2018 GLOF-masked annual maxima belong to the same
        population as the pre-GLOF reference period.
  \item[$H_1^{(2)}$:] The 2011--2018 GLOF-masked annual maxima differ from the pre-GLOF
        reference distribution (evidence of additional hydroclimatic change beyond GLOFs).
\end{itemize}
Failure to reject $H_0^{(2)}$ at $\alpha = 0.05$ supports the stationarity of the
natural-flood regime and attributes the observed distributional shift entirely to the
direct GLOF signal.
